\section{Method}
\label{sec:method}

\subsection{Offline extraction (zero LLM generation cost throughout)}
\label{sec:method:offline}

\paragraph{Two-pass concept extraction} (Phase 25+35; 2 forwards per chunk):
\begin{enumerate}
\item \textbf{Pass 1, concepts}: concern prompt (\emph{``What concepts does this
text discuss?''}) $\to$ full-layer depth gradient + position $-1$ readout
$\to$ triple filtering (ASCII / prefill blacklist / POS) $\to$ BM25-based BPE
completion $\to$ corpus validation with DF $\geq 2$. Accuracy 80\% (human
review).
\item \textbf{Pass 2, role expansion}: inverted prefill (\emph{``The concepts
are: \{concepts\}''}) $\to$ scan the workspace at concept positions $\to$
document-specific role words (e.g., Nutrition $\to$ Education, 71\%
document-specific). The same forward pass also yields, at zero marginal cost,
each concept's conditional vector $ws$ (the residual after
\texttt{lens.transport}), the vector asset used in \S\ref{sec:math}.
\end{enumerate}

\paragraph{Text-side entity detection} (model-free): capitalized spans +
terminology spans + frequency/rule filtering. Phase 52 proves the J-Lens
cannot read out low-frequency proper names (a structural limitation of 7B
models), but proper names are explicitly present in the
text---\emph{detection does not require generation}. This yields 11{,}530
(medical) and 16{,}186 (novel) entities with 100\% detection on the failure
list.

\paragraph{Relation extraction} (1 forward per pair): dual-concept concern
prompt (\emph{``The relationship between A and B is''}) $\to$ position $-1$
$\to$ relation words; 100\% coverage (Phase 27), refined with V3 priors +
prefill scan (Phase 37). \textbf{Works for rare entities as well} (Phase 53
S2, 62.5\%)---the model cannot read out the entity token but can read out the
relation word (common vocabulary).

\paragraph{Entity disambiguation} (1 forward per pair): dual-concept concern +
restricted-classification prefill (\emph{``Answer:''}); accuracy 0.80, beating
the bge cosine baseline of 0.60 (Phase 44). Key engineering finding: the
language prior of the prefill wording directly determines the readout
(\emph{``They are''} locks the output to \emph{``related''}); a restricted
answer set is essential.

\paragraph{Prompt design law} (three independent replications: Phase 16/24, 42,
43): \emph{cloze-style prompts read out syntactic continuations rather than
semantic content; J-Lens readout must be anchored at entity positions (concept
words in the prefill) or on a restricted answer set.}

\subsection{Data structures: matrix formalization}
\label{sec:method:matrix}

The extraction products are organized into four matrices/tensors, all amenable
to linear-algebraic operations:
\begin{itemize}
\item $\mathbf{M}$: concept/entity--document map ($N \times N_{\text{chunks}}$,
sparse), $M[i,j] = \mathrm{IDF}(i) \times \mathrm{BM25}_{tf}(i,j)$;
\item $\mathbf{W}$: concept-relation adjacency ($N \times N$, typed,
probability-weighted); third-order form $\mathcal{T}[i,r,j]$ (relation-type
tensor);
\item $\mathbf{E}_{ws}$: concept workspace conditional vectors
($N \times 3584$), produced for free in Pass 2;
\item $\mathbf{R}$: concept-role matrix (optional).
\end{itemize}

\subsection{Online retrieval}
\label{sec:method:online}

\paragraph{Seed + expand} (the only integration pattern verified correct,
Phase 41B/42):
\begin{enumerate}
\item Query $\to$ bge-m3 $\to$ top-$K$ seed chunks (seeds fixed);
\item concept propagation: $q \times \mathbf{M}$ (one matrix multiply);
\item relation expansion: $\mathbf{W} \times q$ (one matrix multiply);
\item expansions fill the remaining top-$K$ slots ($\text{seed}_k = 5$).
\end{enumerate}

\paragraph{Reimplemented graph retrieval algorithms} (\S\ref{sec:replacement}):
\begin{itemize}
\item \textbf{dual-level} (LightRAG): query $\to$ dual-path vector matching
over entities/relations $\to$ one-hop neighbors $\to$ chunk MAX aggregation
interleaved with naive retrieval;
\item \textbf{PPR} (HippoRAG2): $\alpha\,(I - (1-\alpha)\mathbf{P})^{-1} s$,
with phrase seeds (query-to-triple top-5 + recognition-memory filtering) and
passage seeds (bge similarity $\times$ weight factor).
\end{itemize}

\subsection{Final architecture: four heterogeneous layers}
\label{sec:method:arch}

\begin{table}[t]
\centering
\caption{The four-layer heterogeneous architecture.}
\label{tab:arch}
\begin{tabular}{@{}llll@{}}
\toprule
Layer & Content & Source & Cost \\
\midrule
Symbolic/abstract & Concepts, relations, roles, disambiguation & J-Lens workspace readout & 1 forward per chunk(/pair) \\
Surface & Named entities, proper names, terms & Text-side detection (rules/BM25) & $\sim 0$ \\
Geometric & Attributes, similarity, grounding validation & $\mathbf{E}_{ws}$ / $\mathbf{M}$ vectors & $\sim 0$ \\
Relational & Multi-hop, link completion & Powers of $\mathbf{W}$ / CP tensor & $\sim 0$ \\
\bottomrule
\end{tabular}
\end{table}

Division-of-labor principle: \emph{the model does abstraction (the most
expensive and error-prone part of generation), deterministic methods do the
surface (the part that was free anyway), and escape paths are reserved for the
hard cases} (\S\ref{sec:escape}). Layers align through concept symbols but do
not share a vector space (the decoupling verdict of
\S\ref{sec:math:decoupling}).
